% introduction_two_structures.tex
\subsection{Two Relational Structures Over the Same Items}

Psychology's statistical machinery for such relational webs predates construct-validity vocabulary by half a century. Holding character ``a definite and durable `something'\,'' and so measurable, \citet[p.~181]{galton-1884} mapped it from the dictionary, not behavior, estimating in Roget's Thesaurus ``fully one thousand words expressive of character,'' distinct yet overlapping in meaning---founding the lexical tradition and anticipating the semantic-similarity analyses below. \citet{spearman-1904} supplied the inferential engine of ``correlational psychology'': diverse intellectual performances correlate positively; their common variance marks a single general function, estimable with its per-measure saturation from the intercorrelation matrix. Factor analysis---low-dimensional structure from a high-dimensional relational matrix---descends from this move; \citet{joreskog-1969} made it confirmatory---any subset of loadings, factor correlations, and uniquenesses fixed in advance, the rest maximum-likelihood estimated with a likelihood-ratio test of fit---and measurement models testable hypotheses. \citet{goldberg-1990} rejoined the threads: virtually identical five-factor structures emerged across 1{,}431 trait adjectives, ten extraction--rotation combinations, and self- and peer ratings, establishing the Big Five as a robust summary of the English trait lexicon. The object was always a relational matrix whose variables, in the lexical tradition, were language; semantic factor analysis changes the source, not the logic.

A construct enters empirical practice through its items---the researcher's most concrete public commitment, in ordinary language, to the experiences, behaviors, and situations it involves: for \citet{messick-1995}, the content aspect of validity. Reflective or formative \citep{bollen-lennox-1991}, items encode the intended meaning more exactly than any definition. If, as \citet[\S43]{wittgenstein-1953} held for a large class of cases, ``the meaning of a word is its use in the language,'' that meaning is public and analyzable---regularities of use across the language community, prior to and independent of any respondent sample; and a concept needs no sharp boundaries or common essence to do scientific work, its instances sharing only family resemblances \citep[\S\S66--67]{wittgenstein-1953}---apt for constructs whose items cohere without one shared feature. Two relational structures over the same items follow: items cluster behaviorally in the person-by-item response matrix because the same people endorse them, and semantically in the language itself because their words are used alike---recovered, respectively, by classical and semantic factor analysis. Nor is wording behaviorally peripheral: \citet{clark-watson-2019} observe that exact phrasing can greatly influence the construct measured---almost any negative mood term virtually guarantees a substantial negative-affectivity component---and that a single item carries interpretable variance at several construct-hierarchy levels; a method analyzing wording directly is thus a missing instrument.
